% !TEX root = ../main.tex
\chapter{서론과 연구 문제}
\label{ch:introduction}

디파이(DeFi, 탈중앙 금융)는 누구나 볼 수 있는 블록체인 위에서 돈(자본)이 쓰일 곳을 정하는 방식을 바꾸었어요. 이제는 예전의 은행 같은 중간 기관 없이도, 스마트 계약(블록체인 위에서 저절로 실행되는 프로그램)으로 돈을 빌려주고, 빌리고, 빌린 돈을 보태 거래할 수 있어요 \parencite{schar2021defi}. 이런 서비스를 대부분 쓰려면, 토큰 주인(허락을 해 주는 쪽)이 먼저 계약에 자기 토큰을 옮길 권한을 내줘야 해요. ERC-20(토큰 공통 규칙) 표준은 이렇게 내준 권한을 하나하나 허락(allowance, 토큰을 얼마까지 대신 써도 된다는 허락)으로 기록해요 \parencite{vogelsteller2015erc20}. 그리고 사용자도, 사용자가 서명할 때 쓰는 프로그램도, 다른 서비스(프로토콜)도 모두 어떤 스펜더 계약(허락을 받아 대신 쓰는 계약)에게 허락을 줄지 정해야 해요. 그런데 진짜 이름 대신 주소로만 활동하는 세상(가명성)에서는 이 결정을 잘 내리기가 어려워요. 이 논문의 가장 중요한 문제는, 앞서 내준 권한들의 공개 기록이 한 계약에 대해 무엇을 말해 줄 수 있는지예요.

이 논문은 ERC-20 허락 화살표(엣지)가 스마트 계약의 평판 점수로 쓰는 페이지랭크(PageRank)에 무엇을 보태는지 물어요. 허락 화살표는 \texttt{approve}(허락하기)와 \texttt{permit}(서명 허락)이 남기는 \texttt{Approval} 이벤트(계약이 남기는 기록)에서 나와요. 엔도스랭크(EndorseRank)는 이 화살표 위에서 계산한 페이지랭크예요. 화살표 무게(가중치)는 \textcite{do2023awp}의 AWP(적응형 가중 페이지랭크)에서 가져와요. AWP는 송금(transfer)을 바탕으로 한 방법이에요. 이 무게는 널리 거래되는 토큰 다섯 가지의 금액을 달러(USD)로 나타낸 값에 매겨요. 두 점수는 아비트럼 원(Arbitrum One)의 같은 계약들에 대해 계산해요. 그 곁에서 두 종류의 화살표를 함께 걷는 결합 걷기(화살표를 따라 점에서 점으로 옮겨 다니는 계산)와, 이 걷기들이 매끄럽게 다듬는(평활하는) 원시 차수(화살표를 그냥 센 수)도 함께 계산해요. 그리고 AWP를 검사할 때 쓴 순위 상관 방법(두 순위가 얼마나 비슷한지 재는 방법)으로 점수들을 점검해요 \parencite{do2023awp}. 이 연구에서는 그 방법에 짝지은 대비(같은 대상에서 두 점수를 맞대어 보는 비교), 원시 차수 기준선(비교 기준), 시간 홀드아웃(나중 기간을 떼어 두고 맞혀 보는 시험) 두 번을 더했어요. 두 번째 홀드아웃은 그 기간의 로그(기록)를 뽑기 전에 미리 등록(계획을 먼저 공개해 두는 일)해 둔 기간에서 했어요. 주된 분석의 금액은 등록 분석(미리 등록한 그대로 한 분석)의 결과를 알고 난 뒤에 고쳤어요(\ref{sec:supplementary-extension}절). 등록 분석은 부록에 남겨 두었어요(부록~\ref{sec:map-registered}). 점수들은 온체인(블록체인 위의) 이벤트를 써요. 그리고 금액의 값을 매길 때는 누구나 볼 수 있는 하루 토큰 가격도 써요. 온체인 평판 점수가 무엇인지는 다른 주요 용어들과 함께 \ref{sec:research-problem}절에서 뜻을 밝혀요.

\dissertationSubheading[sec:background]{배경: 디파이, 담보, 가명성}

누구나 볼 수 있는 블록체인은 중앙 기관 없이도 프로그램으로 거래를 마무리(결제)할 수 있게 해 줘요 \parencite{nakamoto2008bitcoin}. 스마트 계약은 약속의 조건을 코드로 적어 둘 수 있어요 \parencite{szabo1997formalizing}. 이더리움(Ethereum)에서는 이런 계약이 자기 코드의 규칙에 따라 토큰을 가지고 있거나 옮길 수 있어요 \parencite{buterin2014ethereum}. 이런 계약으로 만든 디파이 플랫폼은 돈을 맡아 주는 기관(수탁 기관)을 거치지 않고 블록체인 위에서 거래를 마무리해요. 그리고 서비스들은 이 계약들을 쌓아 올려, 빌려주기(대출), 자동 마켓 메이킹(프로그램이 자동으로 사고팔아 주는 일), 무기한 선물 거래를 하는 탈중앙 시장을 만들어요 \parencite{werner2022sok}. 아비트럼 원과 다른 레이어 2(L2) 네트워크에서는 수수료가 충분히 낮아서, 보통 사용자도 자주 거래해요. 그리고 모든 송금, 허락, 서비스 이벤트를 주소 하나하나마다 살펴볼 수 있어요.

이렇게 다 보이는데도, 디파이의 신용 시장은 여전히 과담보(빌리는 돈보다 담보를 더 맡기는) 상태예요. 대부분의 서비스는 빌리는 사람이나 트레이더(거래하는 사람)에게 포지션(거래로 걸어 둔 몫)의 값보다 더 많은 담보를 맡기게 해요. 거래 상대가 가명의 계정이고, 그 계정의 주인이 누구인지 확인되지 않았기 때문이에요 \parencite{harvey2021defi}. 가명성은 누구나 볼 수 있는 장부의 기본 성질이에요. 누구나 거의 돈을 들이지 않고 주소를 만들 수 있어요. 그리고 한 사람이나 한 단체가 주소를 많이 거느릴 수 있어요 \parencite{douceur2002sybil}. 보통의 신용 등급은 개인의 지난 기록에 기대는데, 그런 기록은 이런 환경에 잘 옮겨지지 않아요 \parencite{packin2024credit}. 그래서 서비스들은 한 사람 한 사람의 평판 대신, 쌓아 둔 담보, 청산(손실이 너무 커져 거래가 강제로 끝나는 일) 장치, 서비스 규칙으로 위험을 다스려요. 연구자들은 오프체인(블록체인 밖의) 신용 기록 대신 쓸 것으로 온체인 신호에 눈을 돌렸어요. 그런 연구로는 송금 그래프(점과 화살표로 이은 그림)의 중심성(그래프에서 얼마나 가운데에 있는지) \parencite{do2023hodgerank}, 네트워크를 따라 위험이 퍼지는 것(위험 전파)을 보는 방법 \parencite{lin2025riskprop}, 전망 이론에 바탕을 둔 신용 점수 매기기 \parencite{hassija2020lending}, 대출 서비스를 대상으로 한 기계 학습 처리 과정(파이프라인) \parencite{kotb2023credit}이 있어요. 이 방법들은 모두 빌리는 사람이나 그 뒤에 있는 계정에 점수를 매겨요.

담보는 대출 서비스를 빌리는 사람들로부터 지켜 줘요. 하지만 어떤 계약이 자기 토큰을 써도 되는지 정해야 하는 주인에게는 아무 도움이 되지 않아요. 이 결정에서는 가명의 거래 상대가 바로 그 계약이에요. 라우터(거래를 이어 주는 계약), 금고(vault, 토큰을 맡아 두는 계약), 애그리게이터(여러 거래소를 모아 이어 주는 계약), 대출 풀(빌려줄 돈을 모아 둔 계약)이 주인의 토큰을 옮기려면, 먼저 주인이 그 계약에 허락을 내줘야 해요. 이 허락은 다 쓰이거나, 새 허락으로 바뀌거나, 거두어질(허락 취소) 때까지 계속 효력이 있어요 \parencite{vogelsteller2015erc20}. 이더리움에서 허락은 대부분 금액에 끝이 없는 무제한 허락이에요 \parencite{wang2022unlimited}. 그래서 나쁜 계약이나 나중에 공격자 손에 넘어간 계약에 허락을 내주면, 주인이 가진 그 토큰 전부가 위험해져요. 피싱 사이트(진짜처럼 꾸민 가짜 사이트)는 이 점을 노려요. 사용자를 꾀어 자기 자산을 공격자에게 넘기는 거래에 서명하게 만들어요 \parencite{he2023txphishscope}. 이런 일을 하는 도구 꾸러미는 이제 서비스처럼 팔리기까지 해요 \parencite{he2025drainer}. 서비스도 자기 계약이 다른 계약을 허락할 때 같은 결정을 해야 해요. 금고가 대출 풀에 넣어 주는 토큰을 그 대출 풀이 끌어가게 허락할 때가 그런 예예요. 계약의 코드는 공개되어 있어요. 하지만 누가 그 계약을 올렸는지(배포했는지), 누가 그 계약을 바꿀 수 있는지, 그 계약이 허락을 어떻게 쓸지는 알 수 없을 때가 많아요.

실제로 스펜더 계약에 대한 믿음은 몇 가지 방법으로 생겨요. 주인은 보통 지갑 프로그램이 서명하라고 허락을 보여 주는 바로 그 순간에 결정해요. 그런데 지갑이 토큰 허락을 보여 주고 확인하는 방식에 약점이 있어서, 그 순간을 노리기가 쉬워요 \parencite{adamczyk2025approval}. 코드 분석은 계약을 그 안의 함수로 판단해요. 예를 들어 러그 풀(만든 사람이 돈을 가지고 달아나는 사기)을 하려고 쓴 함수를 찾아 표시할 수 있어요 \parencite{lin2024crpwarner}. 실제로 쓰이고 있는 평판 도구는 사람 단위로 작동해요. 휴먼 패스포트(Human Passport)는 온체인과 오프체인의 자격 증명을 무게를 매겨 모아요. 그리고 그것으로 한 주소를 서로 다른 진짜 사람 한 명이 움직이는지 판단해요. 이 도구는 시빌 공격(가짜 주소를 많이 만들어 속이는 공격)을 막는 문 역할을 해요. 단계별로 줄을 세우는 순위 역할을 하지는 않아요 \parencite{humanpassportnddocs}. 이것은 열쇠(개인 키)를 가진 사람에 대한 질문이에요. 계약에는 이 질문과 바로 짝이 되는 질문이 없어요. 계약이 하는 일은 그 코드와, 그 계약을 고칠(업그레이드할) 수 있는 누군가가 정하기 때문이에요. 이 방법들 가운데 어느 것도 다른 주인들이 내린 결정의 공개 기록을 쓰지 않아요. 곧, 많은 주인이 어떤 계약을 허락했는지, 그리고 그 주인들도 다시 얼마나 강하게 보증(믿고 맡김)을 받는지를 쓰지 않아요.

그래프 순위 알고리즘은 이런 기록에 잘 맞아 보여요. 페이지랭크와 그 변형들은 되풀이되는 보증으로 점(노드)에 순위를 매겨요. 그래서 한 점의 점수는 다른 점들에게서 받는 보증에 달려 있고, 보증해 주는 그 점들도 다시 보증을 받은 점들이에요 \parencite{page1999pagerank}. 아이겐트러스트(EigenTrust) \parencite{kamvar2003eigentrust}와 트러스트랭크(TrustRank) \parencite{gyongyi2004trustrank} 같은 웹 신뢰 알고리즘은 비슷한 생각을 컴퓨터끼리 직접 주고받는 망(피어 투 피어)의 신뢰와 웹 스팸 문제에 써요. \textcite{do2023hodgerank}에서는 송금 그래프에 페이지랭크와 호지랭크(HodgeRank)를 함께 써서 이더리움 사용자가 얼마나 믿을 만한지 판단해요. 그리고 그 결과를 일종의 사회적 신용이라고 불러요. AWP \parencite{do2023awp}도 송금 그래프로 순위를 매기는 또 다른 알고리즘이에요. AWP는 송금마다 무게를 매기는데, 금액에 대한 상한이 있는 함수와 송금의 나이(얼마나 오래되었는지)에 대한 로지스틱 함수(오래된 일일수록 가볍게 세는 함수)를 써요. 그리고 최근에 송금을 보낸 주소에서 걷기를 다시 출발(재시작)해요. AWP를 만든 사람들은 AWP 점수가 들어오는 송금의 수와 금액과 함께 오르내린다는 것을 보여서 AWP를 점검했어요. 이때 들어오는 송금의 수와 금액은 경제적 중요도의 대리 지표(대신 재는 값)로 보았어요. 이 연구들 때문에 송금 그래프 순위는 기준선이 되었어요. 다른 종류의 화살표로 만든 점수는 이 기준선과 비교해야 해요. 하지만 이 연구들은 주소가 받는 돈으로 주소에 순위를 매겨요. 그리고 쓰는 권한을 분명하게 내주는 일은 빼놓아요. 이 일은 송금과 따로 일어나는 온체인 행동이고, ERC-20 허락으로 기록돼요 \parencite{vogelsteller2015erc20}.

\dissertationSubheading[sec:motivation]{온체인 평판이 필요한 까닭}

스펜더 계약에 평판 점수가 필요한 까닭은 네 가지예요.
\begin{itemize}
\item 허락 결정: 주인, 주인이 서명할 때 쓰는 프로그램, 다른 계약을 끌어와 붙여 쓰는 서비스는 모두 어떤 스펜더를 허락할지 정해요. 다른 주인들이 한 계약을 어떻게 대했는지 순서로 알려 주는 신호가 있으면 이런 결정에 도움이 될 수 있어요. 예를 들어 허락에 서명하기 전에 보여 주는 경고에 넣는 정보 하나로 쓸 수 있어요. 또는 한 서비스가 자기와 연결해도 된다고 허락하는 계약 목록을 만들 때 넣는 정보 하나로 쓸 수 있어요.
\item 검사할 수 있음: 사람이 골라 만든 목록과 오프체인 점수는 속을 들여다볼 수 없어요. 하지만 공개된 이벤트 자료와 공개된 토큰 가격으로 계산한 그래프 점수는, 같은 입력을 가진 사람이라면 누구나 다시 계산하고 검사할 수 있어요. 이것은 탈중앙 금융에서 속을 알 수 없는(블랙박스) 신용 점수를 걱정하는 목소리에 어느 정도 답이 돼요 \parencite{packin2024credit}.
\item 운영 비용: 예를 들어 허락을 묻는 창이 뜰 때마다 새로 계산하는 것처럼 자주 새로 고치는 점수는, 다시 계산하는 데 비용이 적게 들어야 해요.
\item 표준화: 다른 시장에서는 거래 상대를 가려내려고 심사 비율(상대를 걸러 내는 숫자 기준)을 써요. 하지만 디파이에는 계약 단위에서 이 일을 할, 순서를 나타내는 표준 언어가 없어요. 이 논문에서 살펴본 점수들이 어디에서 그 역할을 할 수 있고 어디에서는 할 수 없는지는 \ref{ch:discussion}장에서 이야기해요.
\end{itemize}

송금에 바탕을 둔 AWP는 공개되어 있고 검사할 수도 있어요. 하지만 AWP의 그래프는 돈이 오갈 때마다 커져요 \parencite{do2023awp}. 허락 그래프는 짜임새가 다른 또 하나의 길이에요. 주인이 \texttt{approve}를 부르거나 \texttt{permit}에 서명할 때마다 \parencite{lundfall2020permit}, 토큰 계약은 스펜더가 그 토큰을 얼마까지 옮길 수 있는지 기록해요 \parencite{vogelsteller2015erc20}. 그리고 새 허락은 예전 허락을 덮어써요. 그래서 주인--스펜더 쌍과 토큰마다 마지막 허락(가장 최근 허락) 하나가 그 쌍의 지난 기록을 대신해요. 허락은 송금보다 드물기도 해요. 그래서 같은 계약들에서 보증 그래프 $G_{\text{approve}}$의 화살표 수는 송금 그래프 $G_{\text{transfer}}$의 절반도 안 돼요. 페이지랭크를 푸는 계산은 반복(한 바퀴)마다 드는 비용이 화살표 수 $|E|$에 정비례해서 커져요 \parencite{leskovec2020mining}. 그래서 그만큼 시간이 덜 걸려요(\ref{sec:benchmark-results}절). 두 경우 모두 같은 풀이 프로그램(solver)이 돌아가요. 그래서 이 차이는 더 듬성듬성한(화살표가 적은) 그래프에서 오는 것이고, 알고리즘이 나아진 것은 아니에요. 주장~C1은 이 차이를 그래프를 나타내는 방식의 성질로 다뤄요.

효율 말고도, 허락은 뜻도 달라요. 송금은 다른 곳으로 가는 길에 거쳐 가는 것일 수도 있고, 마켓 메이킹일 수도 있고, 그냥 가만히 받은 것일 수도 있어요. 반면 허락은 돈을 써도 된다는 권한을 일부러 내주는 일이에요. 페이지랭크는 화살표를 인용(다른 글을 가리켜 추천하는 것)처럼 읽어요 \parencite{page1999pagerank}. 허락을 보증 화살표로 다루면, 이 읽는 법을 주인과 서비스가 이미 기대고 있는 믿음 관계로 옮겨 올 수 있어요. 그리고 보증을 받는 점은 스펜더 계약, 곧 주인이 결정을 내려야 하는 바로 그 대상이에요. 이렇게 하면 순위가 달라지는지, 그리고 그 순위가 허락 점검, 송금 점검, 나중 활동과 맞는 정도가 달라지는지는 아비트럼 원의 매칭 코호트(모든 점수가 똑같이 점수를 매기는 계약 묶음)에 든 계약 $n=15{,}107$개로 살펴봐요.\footnote{매칭되었다는 말은 모든 방법이 같은 계약들에 점수를 매긴다는 뜻이에요. 이 계약 목록은 송금은 하나도 보지 않고, 토큰 다섯 가지의 허락 로그만으로 뽑았어요. 뽑는 규칙은 \ref{sec:data-sources}절에 있어요.}

\dissertationSubheading[sec:research-problem]{연구 문제}

AWP는 송금 기록에서 무게를 단 합과 로지스틱 감쇠를 써서 믿음을 짐작해요 \parencite{do2023awp}. AWP는 송금마다 경제적인 무게를 줘요. 하지만 AWP는 (i) 송금 기록 전체를 모아 계산해야 하고, (ii) 값진 주소를 큰 금액의 송금을 많이 받는 주소로 보며, (iii) 쓰는 권한을 내주는 일에는 아무 역할도 주지 않아요. 이 논문이 살펴보는 점수인 엔도스랭크는 AWP의 화살표 무게와 풀이 프로그램을 그대로 써요. 하지만 방향이 있는 그래프는 주인--스펜더 쌍과 토큰마다의 마지막 허락으로 만들고, 그 금액은 달러로 매겨요. 그렇게 만든 그래프는 더 듬성듬성하고, 다른 신호를 담고, 각 관계에서 다른 쪽에 순위를 매겨요. 허락은 주인에서 스펜더 쪽을 가리켜요. 그래서 엔도스랭크는 라우터, 금고, 대출 풀처럼 허락을 받는 계약에 점수를 매겨요. 그리고 각 주인은 자기가 허락한 스펜더들에게 자기 보증을 넘겨줘요. 주인은 외부 소유 계정(EOA, 개인 키로 움직이는 계정)일 수도 있고 다른 계약일 수도 있어요.

주된 연구 질문은 이거예요. 아비트럼 원의 스마트 계약에 매기는 평판 점수로서, ERC-20 허락 화살표는 페이지랭크에 무엇을 보탤까요? 송금 걷기와 비교하면, 그리고 두 걷기가 매끄럽게 다듬는 원시 차수와 비교하면 어떨까요? 연구 질문(RQ) RQ1--RQ5는 이 질문을 잘게 나눈 거예요(\ref{sec:research-questions}절). 두 종류의 화살표는 같은 주소들 위에 있어요. 그래서 걷기 하나가 미리 정해 둔 섞는 비율 $\lambda$ 하나로, 점마다 두 화살표를 섞을 수도 있어요. 이렇게 정한 점수가 C-PR(결합 페이지랭크, 두 그래프를 함께 걷는 점수)이에요. $\lambda=1$과 $\lambda=0$에서는 C-PR이 한 층(그래프 하나)만 걸어요. 두 종류의 화살표는 같은 계약들에 대해 서로 다른 행동을 기록해요. 그래서 이렇게 섞은 것을 시험해 볼 가치가 있어요. 이미 송금 걷기를 돌리고 있는 운영자는 점수를 하나 더 돌리는 대신, 그 걷기에 허락 화살표를 더할 수 있어요. 아직 풀리지 않은 질문은, 그렇게 했을 때 걷기가 송금 신호를 잃지 않으면서 허락 신호도 담을 수 있느냐예요.

점수가 자기 화살표의 통계와 잘 맞거나, 같은 화살표가 앞으로 더 생길 것을 예측(미리 맞히기)한다고 해도, 그것만으로는 그 점수가 알맹이 있는 뜻에서 평판을 잰다고 아직 말할 수 없어요. 그래서 연구 질문을 새로 더하지는 않는 비교를 하나 더 해요. 이 비교는 두 종류의 화살표 어느 것으로도 만들지 않은 결과를 써요. 이 결과는 동결(점수를 고정함) 뒤에 관측했어요. 바로 라벨 기간(나중에 채점에 쓰는 정답 값을 모으는 기간)에 거래한 GMX~V2(무기한 선물 거래 서비스) 계정 가운데 계약인 계정들이, 그 뒤에 청산을 얼마나 피했는지예요(\ref{sub:neutral-label-protocol}절). 이 비교의 규칙은 두 번째 홀드아웃과 함께 등록했어요. 그리고 고친 점수에 재분석으로 적용해요. 이 몇 달 동안 이런 계약은 몇 개밖에 거래하지 않았어요. 그래서 이 비교는 작은 표본에 기대고 있어요.

점수들은 표~\ref{tab:primary-methods-preview}에 있어요. 표~\ref{tab:hrq-claim-map}에는 가설(H1--H3), 연구 질문(RQ1--RQ5), 그리고 \ref{ch:results}장에서 나온 실험 결과 주장(C1--C6)을 정리했어요.

\begin{table}[htbp]
\centering
\caption{이 논문에서 평가한 점수들(매칭 계약 코호트, $n=15{,}107$).}
\label{tab:primary-methods-preview}
\fitwidth{%
\begin{tabular}{p{0.18\linewidth}p{0.28\linewidth}p{0.48\linewidth}}
\toprule
점수 & 입력 그래프 & 역할 \\
\midrule
EndorseRank & 다섯 토큰의 마지막 ERC-20 허락($G_{\text{approve}}$), 달러 금액, AWP처럼 무게를 매김; 고르게 다시 출발(균등 재시작) & 허락으로 맡김; 이 연구에서 살펴보는 대상. AWP처럼 활동에 비례해 다시 출발하는 변형도 옆에 함께 보고해요 \\
AWP & 다섯 토큰의 송금($G_{\text{transfer}}$), 달러 금액 & 송금 흐름; \textcite{do2023awp}의 발표된 송금 그래프 방법. 금액 무게의 기울기는 금액의 중앙값으로 정해요(\ref{sub:awp-graph}절) \\
C-PR & 두 층 모두, 달러 금액(허락 금액은 상한으로 자름), 점마다 무게 $\lambda$로 섞음 & 두 종류의 화살표가 합쳐지는지 시험하는 데 쓰는 결합 연산자(RQ5). $\lambda=1$과 $\lambda=0$에서는 한 층만 걸어요. 이 점수의 빼 보기 실험(한 부분을 빼고 다시 해 보는 실험)은 S-PR(보증 씨앗 페이지랭크)이에요. S-PR은 허락 걷기에서 다시 출발하는 송금 걷기예요 \\
원시 차수 & $G_{\text{approve}}$와 $G_{\text{transfer}}$, 점마다 셈 & 받은 허락 수(허락 입차수)와 송금 입차수(송금을 보내온 서로 다른 주소 수). 페이지랭크들이 매끄럽게 다듬는 그 점 둘레의 수예요. 같은 기간에서는 대리 지표, 기간 밖에서는 기준선으로 써요 \\
\bottomrule
\end{tabular}
}
\end{table}

주요 용어의 뜻은 다음과 같아요.

\begin{itemize}
\item \emph{디파이(DeFi):} 빌려주기, 거래, 빌리기 같은 금융 서비스예요. 누구나 볼 수 있는 블록체인 위에 스마트 계약으로 만들어져 있어요. 그리고 대개 돈을 맡아 두거나 거래를 마무리하는 일을 책임지는 중앙 운영자 없이 돌아가요 \parencite{schar2021defi}.
\item \emph{스마트 계약:} 블록체인 위에서 스스로 실행되는 프로그램 코드예요. 믿고 맡길 중간 기관 없이 약속을 실행해요 \parencite{szabo1997formalizing}. 이 논문에서 허락 로그와 송금 로그는 토큰 다섯 가지, 곧 WETH, USDC, USDT0, ARB, WBTC의 계약이 남겨요. 그리고 점수를 받는 주소들도 계약이에요.
\item \emph{주소, 계약 계정, EOA:} 주소는 계정 하나를 가리켜요. 계약 계정은 그 안에 저장된 코드에 따라 움직여요. 반면 외부 소유 계정은 개인 키로 움직여요 \parencite{antonopoulos2018mastering}. 이 논문에서는 정해 둔 블록 하나에서 \texttt{eth\_getCode}가 바이트코드(계약 프로그램의 기계어)를 돌려주면 그 주소를 계약으로 쳐요. 바이트코드가 없거나, EIP-7702 위임 표시 \parencite{buterin2024eip7702}가 있으면 그 주소는 EOA예요(\ref{sec:data-sources}절). EOA는 그래프에 주인, 보낸 쪽(송신자), 받는 쪽(수신자)으로 들어오고, 점수를 받는 일은 결코 없어요. 오프체인 신원은 미리 가정하지도 않고 알아내지도 않아요.
\item \emph{온체인 평판 점수:} 온체인 자료로 계산한 수예요. 한 주소가 자기와 주고받는 주소들 사이에서 상대적으로 어디쯤에 있는지 나타내요 \parencite{do2023hodgerank}. 이 논문에서 점수를 받는 주소는 계약이에요. 오프체인에서 가져오는 입력은 토큰 다섯 가지의 목록과, 그 금액의 값을 매기는 하루 달러 가격뿐이에요(\ref{sec:amount-revision}절). 점수는 순서를 나타내는 값이고, 순위를 매기려고 써요. 점수 크기 자체에는 뜻이 없어요.
\item \emph{페이지랭크:} 감쇠 계수(화살표를 따라갈 확률)를 써서 방향이 있는 그래프 전체로 영향력을 퍼뜨리는 그래프 순위 알고리즘이에요 \parencite{page1999pagerank}. AWP는 송금 그래프 위의 페이지랭크예요. 화살표마다 송금의 금액과 오래된 정도로 무게를 매기고, 최근에 활발했던 보낸 쪽에서 다시 출발해요 \parencite{do2023awp}. 이 논문에서 금액은 달러로 나타내요. 금액 무게는 금액이 중앙값일 때 2분의 1이고, 금액이 클수록 1에 가까워져요. 엔도스랭크는 AWP의 화살표 무게를 마지막 허락 그래프에 쓰고, 고르게 다시 출발해요(\ref{sec:graph-construction}절).
\item \emph{결합 페이지랭크(C-PR):} 허락 층과 송금 층을 함께 걷는 무작위 걷기 하나예요. 각 점에서는 그 점이 들어 있는 층들에 대해, 두 층의 행(이동 행렬에서 그 점의 줄)을 무게 $\lambda$와 $1-\lambda$로 섞어요(\ref{sub:coupled-graph}절). 보증 씨앗 페이지랭크(S-PR)는 송금 걷기를 그대로 두지만, 허락 걷기 결과에서 다시 출발해요. 트러스트랭크가 씨앗 집합(믿을 만하다고 미리 고른 점들)에서 다시 출발하는 것과 같아요 \parencite{gyongyi2004trustrank}.
\item \emph{ERC-20 허락 / Approve / Permit:} 스펜더가 주인 대신 주인의 토큰을 얼마나 쓸 수 있는지 기록하는 계약 함수예요 \parencite{vogelsteller2015erc20}. \texttt{Approve}는 보통의 온체인 호출이에요. EIP-2612 \texttt{permit}은 서명으로 같은 상태 변화를 만들어요 \parencite{lundfall2020permit}. 엔도스랭크는 주인--스펜더 쌍마다, 그리고 다섯 토큰마다 마지막 허락을 써요. 이 허락은 기준일(값을 매기는 기준이 되는 날) 그날의 토큰 가격으로 값을 매기고, 0보다 클 때만 남겨요.
\item \emph{보증 그래프:} 방향이 있는 그래프인데, 화살표가 실제로 자산을 옮긴 것이 아니라 쓰는 권한을 분명하게 맡긴 것을 나타내요. 화살표 $u \to v$는, 주인 $u$가 $v$에게 $u$의 토큰을 써도 된다고 허락했다는 뜻이에요 \parencite{vogelsteller2015erc20}. 이 주인은 EOA일 수도 있고 계약일 수도 있어요. 점수를 받는 스펜더 $v$는 계약이에요.
\item \emph{분야 일치도:} 평판 순위가 눈으로 볼 수 있는 대리 지표의 순위와 얼마나 맞는지 나타내는 정도예요. 스피어먼의 로($\rho$) \parencite{spearman1904proof}와 켄달의 타우($\tau$, 두 줄 세우기가 얼마나 같은 순서인지 나타내는 수) \parencite{kendall1938rank}로 재요. 점수 자신의 화살표로 만든 대리 지표와 맞는 것은, \textcite{cronbach1955construct}에서 말하는 뜻의 의도한 구성개념 점검(점수가 재려고 하는 생각을 정말 재는지 보는 점검)으로 읽어요. 어떤 방법은 다른 점검보다 자기 구성개념의 그래프와 더 잘 맞을 수 있어요.
\end{itemize}

\dissertationSubheading[sec:supplementary-extension]{시간 홀드아웃}

같은 기간의 허락 계수(표~\ref{tab:alignment-hybrid})는 의도한 구성개념 점검이에요. $G_{\text{approve}}$ 위의 페이지랭크는 각 스펜더의 받은 허락 수와 잘 맞아야 해요. 페이지랭크는 그 수를 그래프 전체에 걸쳐 매끄럽게 다듬기 때문이에요. 점수가 같은 기간의 그래프 너머의 정보를 전하는지 시험하려고, 점수를 $t_1$(2026년 3월 31일, 23:59:59~UTC)에 동결해요. 그리고 2026년 4월 1일부터 6월 30일까지 센 두 라벨과 견주어요. 하나는 스펜더 계약이 얻는 새 주인--스펜더 허락 쌍(새로 생긴 주인과 스펜더 사이의 허락)이에요(보증 구성개념). 다른 하나는 그 계약의 새 송금자(새로 송금을 보내온 주소)예요(흐름 구성개념). 이 첫 번째 기간인 W0은 스펜더 코호트로 평가해요. 스펜더 코호트는 뽑아낸 허락 자료에서 $t_1$에 다섯 토큰 가운데 하나의 0보다 큰 마지막 허락을 적어도 하나 가진 모든 계약이에요($n=17{,}050$). 또 각 페이지랭크를 그것이 $t_1$에 매끄럽게 다듬는 원시 차수와도 견주어요. 그래서 퍼뜨린 점수가 그 점에서 그냥 센 수보다 얼마나 나은지 바로 볼 수 있어요. 이 차수 대비는 미리 정한 규칙 어디에도 들지 않아요. 그리고 라벨을 본 뒤에 계산했어요. W0에는 규칙 A의 주된 부분이 들어 있어요. 규칙 A는 분석 계획에서 결합 연산자에 대해 정한 규칙이에요. W0에는 같은 몇 달 동안의 GMX~V2 계약 트레이더(GMX V2에서 거래하는 계약)에 대한 탐색적(미리 정하지 않고 살펴본) 실행도 들어 있어요(\ref{sec:holdout-protocol}절).

두 번째 기간인 W1은 2026년 10월 9일 16:49~UTC에 등록했어요. W0의 결과를 안 뒤였고, 2026년 6월 30일 뒤의 날짜가 찍힌 로그를 하나라도 뽑기 전이었어요. 로그를 뽑는 스크립트는 등록 문서와 설정 파일이 커밋되어 있지 않으면 돌아가지 않아요. 그리고 이 스크립트는 두 파일의 해시(파일 내용으로 만든 확인용 값)를 기록해요. W1은 관측 기간(자료를 지켜본 기간)의 끝(2026년 6월 30일, 23:59:59~UTC)에 점수를 동결해요. 그래서 점수는 같은 기간 평가와 같은 그래프에서 나와요. W1의 라벨은 2026년 7월 1일부터 9월 30일까지 세요. W1은 그 동결일에 0보다 큰 마지막 허락을 가진 스펜더 계약에서 두 스펜더 라벨을 다시 재요. 수정 분석(등록 뒤에 금액을 달러로 바꿔 다시 한 분석)에서는 그 허락이 다섯 토큰 가운데 하나의 허락이어야 해요($n=18{,}152$). 그리고 W1은 계약 트레이더를 위한 라벨을 하나 더해요. 그 라벨은 닫은 거래 가운데 청산으로 끝나지 않은 몫이에요. 대상은 계약인 GMX~V2 계정 가운데, 동결일 그래프에 나오고 그 기간에 포지션을 적어도 세 번 닫은 계정이에요($n=70$). 등록 문서는 결합 연산자에 대한 규칙 B를 정해요. 여기에는 규칙 B의 허용 폭, 판정 밖에서 덧붙여 보고하는 대비 세 개, 민감도 분석(조건을 바꿔 결과가 얼마나 움직이는지 보기) 두 개가 함께 들어 있어요. 또 등록 문서는 그 트레이더 라벨에서 허락 쪽 점수와 송금 쪽 점수를 견주는 비교도 정해요. 나올 수 있는 결론마다 어떤 말로 쓸지도 함께 정해요(\ref{sub:fresh-holdout}절과 \ref{sub:neutral-label-protocol}절). 그 비교는 W0을 보조 기간으로 써요. 그리고 등록 문서를 커밋했을 때, W0 트레이더의 청산 라벨은 어떤 점수에 대해서도 아직 계산하지 않았어요.

2026년 10월 10일에 금액을 다루는 방법을 고쳤어요(\ref{sec:amount-revision}절). 같은 기간, W0, W1에서 이 등록 분석의 결과를 모두 안 뒤였어요. 계획은 모든 ERC-20 토큰의 원래 단위(소수점 조정을 하지 않은 가장 작은 단위)를 썼어요. 그리고 거의 모든 금액에서 1에 가까운 금액 무게를 썼어요. 그래서 등록 분석에서는 금액이 엔도스랭크에도 AWP에도 거의 아무 역할을 하지 않았어요. 수정 분석은 모든 금액을 그 토큰의 하루 가격으로 달러로 바꾸어 매겨요. 널리 거래되고 가격을 믿을 수 있는 토큰 다섯 개만 남겨요. 금액 무게의 기울기는 중앙값 금액으로 정해요. 그리고 코호트, 그래프, 대리 지표, 스펜더 라벨을 그 토큰들로 다시 만들어요. 이 수정은 고친 점수를 하나라도 계산하기 전에 커밋했어요. 수정 분석이 이 논문의 주된 분석이에요. 규칙 A와 B, 그리고 트레이더 결과에 대한 비교는 정해 둔 기준값과 문구 그대로 수정 분석에 적용해요. 하지만 고친 점수는 등록 분석의 결과를 안 뒤에 정의했어요. 그래서 수정 분석의 판정은 재분석이고, 등록한 시험이 아니에요. 등록 분석도 같은 판정에 이르렀어요. 부록~\ref{sec:map-registered}에 등록 분석의 결과를 실었고, 표~\ref{tab:registered-vs-revised}은 그 결과를 수정 분석의 결과와 나란히 보여 줘요.

\begingroup
\setstretch{1}\small\setlength{\tabcolsep}{3.5pt}\captionsetup{font=normalsize}
\begin{longtable}{p{0.12\linewidth}p{0.46\linewidth}p{0.36\linewidth}}
\caption{가설, 연구 질문, 자료로 확인한 주장의 요약.}
\label{tab:hrq-claim-map}\\
\toprule
항목 & 다루는 것 & 보고한 곳 \\
\midrule
\endfirsthead
\caption[]{(계속)}\\
\toprule
항목 & 다루는 것 & 보고한 곳 \\
\midrule
\endhead
\bottomrule
\endlastfoot
H1 / RQ1 & EndorseRank와 AWP는 같은 계약들을 서로 다른 순서로 세우나요? & \S\ref{sec:rank-divergence}; 방법 사이의 $\tau$와 95\% 신뢰구간(믿을 만한 범위); 표~\ref{tab:tie-sensitivity}; 주장 C3 \\
H2 / RQ2 & 각 점수는 허락 점검, 그리고 송금 점검과 얼마나 잘 맞나요? & 표~\ref{tab:alignment-hybrid}과~\ref{tab:tie-sensitivity}; 주장 C2 \\
H3 / RQ3 & $n$이 같을 때 각 점수를 계산하는 데 얼마나 드나요? & \S\ref{sec:benchmark-results}; 주장 C1 \\
RQ4 & $t_1$에 동결한 점수가 스펜더 계약의 나중 새 허락 쌍과 새 송금자를 예측하나요? 그리고 자기가 다듬는 원시 차수보다 더 잘하는 페이지랭크가 있나요? & \S\ref{sec:holdout-results}; 표~\ref{tab:holdout-spenders}, \ref{tab:holdout-diff}, \ref{tab:holdout-receiving}, \ref{tab:fresh-holdout}와~\ref{tab:fresh-posthoc}; 주장 C4 \\
RQ5 & 결합 연산자 C-PR은 송금만 걷는 걷기보다 나은가요? 그리고 어느 한 층만 걷는 걷기보다 계약에 순위를 더 잘 매기나요? & \S\ref{sec:coupled-results}, \S\ref{sec:holdout-results}과~\S\ref{sub:fresh-results}; 표~\ref{tab:alignment-hybrid}, \ref{tab:tau-diff-hybrid}, \ref{tab:holdout-diff}, \ref{tab:fresh-holdout}와~\ref{tab:fresh-contrasts}; 주장 C5 \\
C1 & 푸는 시간은 각 그래프의 화살표 수를 따라가요. EndorseRank는 더 듬성듬성한 그래프 위에서 돌아서, 한 번 부르는 데 AWP보다 시간이 덜 걸려요. 반면에 송금 층을 걷는 C-PR과 S-PR은 AWP보다 오래 걸려요 & \S\ref{sec:benchmark-results} \\
C2 & 같은 기간의 일치: 각 점수는 자기가 걷는 화살표 종류의 점검과 가장 잘 맞아요. 이것은 의도한 구성개념 점검이고, 일부는 만든 방식 때문에 저절로 그렇게 돼요. AWP는 들어오는 송금의 금액 합계를 그 개수와 거의 똑같이 잘 따라가요. EndorseRank는 동점을 거의 남기지 않아요. 외톨이 점(화살표가 하나도 없는 점) 규칙은 어떤 계수도 바꾸지 않아요. C-PR 혼합은 두 층 사이에 놓여요 &표~\ref{tab:alignment-hybrid}과~\ref{tab:tie-sensitivity} \\
C3 & EndorseRank와 AWP는 코호트를 서로 다른 순서로 세워요 & 표~\ref{tab:tie-sensitivity}; 그림~\ref{fig:rank-scatter} \\
C4 & 기간 밖에서, 라벨을 본 뒤에 계산한 대비에서: AWP는 두 기간 모두에서 새 송금자를 원시 송금 입차수보다 더 잘 줄 세워요. 반면에 고르게 다시 출발하는 EndorseRank는 두 기간 모두에서 새 허락 쌍을 원시 받은 허락 수보다 못하게 줄 세워요. EndorseRank의 변형 가운데 활동 재시작(활발한 주소에서 더 자주 다시 출발하기)을 쓰는 것은 새 허락 쌍을 차수보다 더 잘 줄 세워요. 하지만 이것은 발견이 아니라 새 자료로 시험할 가설이에요 &표~\ref{tab:holdout-spenders}, \ref{tab:holdout-diff}, \ref{tab:holdout-receiving}와~\ref{tab:fresh-posthoc} \\
C5 & 고친 점수에서, 등록 분석을 위해 정한 규칙으로 한 재분석에서: 두 종류의 화살표를 한 번에 걷는 걷기는 한 층만 걷는 걷기를 이기지 못해요. 규칙 A는 두 부분 모두 실패해요. 규칙 B도 채우지 못해요. 새 허락 쌍에서 송금 걷기보다 나은 점은 다시 나타나요. 하지만 새 송금자에서, 그리고 아슬아슬하게 트레이더의 청산 없이 닫은 비율에서 비열등성(정한 폭보다 더 나쁘지 않음)을 채우지 못하기 때문이에요. AWP와 견주면 뚜렷하게 실패해요. S-PR 빼 보기 실험에서, 허락 화살표는 걷기를 직접 이끌 때만 도움이 돼요. 등록 분석도 같은 판정에 이르렀어요 &표~\ref{tab:alignment-hybrid}, \ref{tab:tau-diff-hybrid}, \ref{tab:holdout-diff}, \ref{tab:fresh-holdout}, \ref{tab:fresh-contrasts}와~\ref{tab:registered-vs-revised}; \S\ref{sub:fresh-results} \\
C6 & 어느 화살표 종류로도 만들지 않은 결과 값에서의 비교, 등록한 규칙을 고친 점수에 적용함: 주된 기간에서 나중에 청산을 피하는 데에 대해, 그냥 센 수의 수준에서는 허락이 송금보다 앞서지 않아요. 페이지랭크 버전들은 어느 쪽이 나은지 가려지지 않아요. 이익을 낸 비율과 손실이 아닌 비율에서는 허락한 주인 수(서로 다른 주인 수)가 송금 수보다 낮아요. 트레이더 표본이 작아요 & 표~\ref{tab:neutral-label-contrasts}, \ref{tab:neutral-label-recipients}와~\ref{tab:registered-vs-revised} \\
\end{longtable}
\endgroup

\dissertationSubheading[sec:research-objectives]{연구 목표}

이 연구에는 목표(O)가 다섯 개 있어요. RQ1--RQ5마다 하나씩이에요. 각 목표는 그것을 실제로 재는 데 쓴 측정과 함께 적었어요. 또 \ref{ch:discussion}장에서 그 목표를 이루었는지 판단할 때 쓰는 기준과 함께 적었어요.

\begin{enumerate}
\item[O1.] 엔도스랭크와 AWP가 같은 코호트의 계약들을 서로 다른 순서로 세우는지 밝혀요(RQ1). 측정: 매칭 코호트에서 두 점수 벡터(점수를 차례로 늘어놓은 것) 사이의 켄달의 $\tau$와 95\% 부트스트랩(다시 뽑기) 구간이에요. 기준: 구간이 $0$(서로 상관없는 순서)도, $1$(완전히 같은 순서)도 품지 않아야 해요.
\item[O2.] 같은 기간에서 각 점수가 각 구성개념과 얼마나 잘 맞는지 평가해요(RQ2). 측정: 매칭 코호트에서 모든 점수와 구성개념에 대해 구한, 지표 가족(비슷한 지표 묶음) 수준의 켄달의 $\tau$와 95\% 구간이에요. 기준: 어떤 점수의 그 가족에 대한 구간이 $0$을 품지 않으면, 그 점수는 그 구성개념과 맞는다고 봐요. 점수가 자기 화살표의 구성개념과 맞는 것은 의도한 구성개념 점검으로 읽어요. 그 대리 지표들이 그래프와 같은 이벤트에서 나오기 때문이에요.
\item[O3.] 같은 계약들에서 각 점수를 계산하는 데 얼마나 드는지 재요. 그리고 뽑은 계약 수에 따라 실행 시간이 어떻게 바뀌는지 재요(RQ3). 측정: 푸는 시간(준비 실행을 한 번 한 뒤, 시간 잰 실행 다섯 번의 평균과 표준편차(SD)), 가장 많이 쓴 메모리, 반복 횟수, $|V|$와 $|E|$예요. 모두 한 가지 절차로 재요. 기준: 시간의 비율은 $|E|$의 비율과 함께만 보고해요. 그래야 어떤 차이든 그래프의 크기 때문인지 따져 볼 수 있어요.
\item[O4.] 라벨 기간 전에 동결한 점수가, 스펜더 계약이 그 기간에 얻는 새 허락 쌍과 새 송금자를 예측하는지 시험해요. 그리고 자기가 다듬는 원시 차수보다 더 잘하는 페이지랭크가 있는지 시험해요(RQ4). 측정: 두 원시 차수를 포함한 모든 점수와 두 라벨 사이의 켄달의 $\tau$와 95\% 구간이에요. 대상은 두 기간의 스펜더 계약이에요. W0에서는 2026년 3월 31일에 점수를 동결하고, 2026년 4월부터 6월까지 라벨을 세요. W1에서는 2026년 6월 30일에 점수를 동결하고, 2026년 7월부터 9월까지 라벨을 세요(\ref{sec:supplementary-extension}절). 각 기간에서 각 페이지랭크를 자기 차수와 견주는 짝지은 대비도 구해요. 기준: 어떤 점수의 구간이 $0$을 품지 않으면, 그 점수는 라벨을 예측한다고 봐요. 페이지랭크는 짝지은 차이의 구간이 $0$보다 위에 있을 때만 자기 차수보다 낫다고 봐요. 차수 대비도, 이 기준도 W0 라벨을 알기 전에 정하지 않았어요.
\item[O5.] 두 종류의 화살표를 가로지르는 한 번의 걷기가 송금만 걷는 걷기보다 나은지 시험해요. 그리고 그 걷기가 어느 한 층만 걷는 걷기를 이기는지 시험해요(RQ5). 측정: $\lambda=0.5$인 C-PR($0.25$와 $0.75$는 민감도를 보는 값)과 S-PR이에요. 이 둘을 한 층만 걷는 두 걷기, 그리고 엔도스랭크와 AWP와 견주어요. 견줄 때는 홀드아웃 라벨과 같은 기간 라벨에서 $\Delta\tau$의 짝지은 구간을 써요. 기준: 각 기준은 등록 분석을 위해, 판정에 쓰는 자료보다 먼저 정했어요. 그리고 등록 분석의 결과를 안 뒤에 같은 기준값으로 고친 점수에 적용했어요. 이것은 등록한 시험이 아니라 재분석이에요. (a) 규칙 A는 로그를 하나도 뽑기 전에, 2026년 10월 9일의 분석 계획에서 정했어요. 규칙 A에 따르면 C-PR은 W0에서 새 허락 쌍에 대해 허락 층만 걷는 걷기보다 못하지 않아야 해요. 또 새 송금자에 대해 송금 층만 걷는 걷기보다 나아야 해요. 그리고 같은 기간에서 송금 가족과 허락 가족에 대해서는 더 나은 층의 구간 안에 있어야 해요. 시빌 안정성에 대해서는 두 층보다 모두 위에 있어야 해요. (b) 규칙 B는 그 뒤 같은 날인 2026년 10월 9일에 등록했어요. W0 결과를 안 뒤였고, W1 로그를 뽑기 전이었어요. 규칙 B에 따르면 C-PR은 W1에서 새 허락 쌍에 대해 송금 층만 걷는 걷기를 이겨야 해요. 그리고 새 송금자와 계약 트레이더의 청산 없이 닫은 비율에 대해서는, 그 걷기보다 허용 폭 $0.01$을 넘게 나쁘지 않아야 해요. 이 허용 폭은 등록 분석에서 그에 맞는 W0 구간의 반폭(구간 너비의 절반)으로 정하고, 올림했어요. 규칙 B와 함께 등록한 민감도 분석은 송금 층 자리에 AWP를 넣고 (b)를 되풀이해요. 이 분석은 규칙을 바꾸지 않아요. C-PR과 S-PR의 다른 모든 대비는 탐색적이에요. 견주는 점수와 라벨을 본 뒤에 계산했기 때문이에요.
\end{enumerate}

\dissertationSubheading[sec:research-questions]{연구 질문}

첫 세 연구 질문에는 가설(H) 세 개가 따라와요. H1: 같은 계약 묶음에서 엔도스랭크와 AWP는 서로 다른 순서를 만들어요. H2: 각 점수는 자기가 걷는 화살표로 만든 점검과 잘 맞아요. H3: $n$의 값이 같을 때, 점수의 계산 시간은 그 그래프의 화살표 수를 따라가요. 그래서 더 듬성듬성한 그래프 위의 엔도스랭크가 AWP보다 빨리 계산돼요. H1--H3은 대부분 만든 방식 때문에 저절로 성립해요. 그리고 이 가설들은 측정 도구를 점검하는 역할을 해요. 틀렸다고 밝혀질 수 있는 시험은 RQ4와 RQ5예요. RQ4의 차수 대비는 등록한 규칙 어디에도 들지 않아요. 그리고 고친 점수에서는 두 기간의 라벨을 모두 안 뒤에 계산했어요. RQ5는 등록 분석을 위해 미리 정한 규칙(\ref{sec:holdout-protocol}절)으로 답해요. 그 규칙을 고친 점수에 재분석으로 적용해요. 표~\ref{tab:hrq-claim-map}에 질문들을 정리했어요.

\begin{enumerate}
\item[RQ1.] 같은 아비트럼 원 계약 묶음에서, 엔도스랭크와 AWP는 서로 다른 순서를 만들까요?
\item[RQ2.] 각 점수는 지금의 허락 점검, 그리고 지금의 송금 점검과 얼마나 잘 맞을까요(스피어먼의 로 $\rho$와 켄달의 타우 $\tau$로 재요)?
\item[RQ3.] 같은 계약 묶음에서 각 점수를 계산하는 데, 시간, 메모리, 반복 횟수로 재면 얼마나 들까요?
\item[RQ4.] 동결일 $t_1$에 계산한 점수가 있을 때, 그 점수는 스펜더 계약이 $t_1$ 뒤에 얻는 새 주인--스펜더 허락 쌍과 새 송금자를 예측할까요? 그리고 $t_1$에 자기가 다듬는 원시 차수보다 그것들을 더 잘 예측하는 페이지랭크가 있을까요?
\item[RQ5.] 허락 층과 송금 층 모두 위에서 도는 결합 페이지랭크는 AWP처럼 송금만 걷는 걷기보다 나을까요? 그리고 같은 기간뿐 아니라 기간 밖에서도, 어느 한 층만 걷는 걷기보다 계약에 순위를 더 잘 매길까요?
\end{enumerate}

RQ4는 같은 두 구성개념의 나중 라벨을 써요. 여기서 순위를 매기는 계약들은 돈을 빌리지 않기 때문이에요. 그래서 이 계약들에는 빚을 갚지 못했는지나 빚을 갚았는지를 나타내는 라벨이 없어요(\ref{sub:non-claims}절). RQ5는 두 화살표 종류를 견주는 데서 나아가, 둘을 합치는 쪽으로 넘어가요. 규칙 A는 C-PR이 어느 한 층만 걷는 걷기를 이기는지 물어요. 규칙 A는 분석 계획에서 $\lambda$와 함께 정했고, W0과 같은 기간에서 판정해요. 송금만 걷는 걷기보다 나은지 보는 규칙은 W0 결과를 안 뒤에 썼어요. 그래서 이 규칙은 W1에서만 판정해요. W1의 로그는 이 규칙을 등록한 뒤에 뽑았어요(\ref{sec:holdout-protocol}절과 \ref{sub:fresh-holdout}절). 두 규칙 모두 등록 분석의 결과를 안 뒤에 고친 점수에 적용해요(\ref{sec:amount-revision}절). 부록~\ref{sec:map-registered}에는 두 규칙의 등록 분석 판정을 실었어요.

\dissertationSubheading[sec:literature-summary]{문헌 검토 요약}

\ref{ch:literature}장은 네 분야를 살펴봐요. 그래프 순위 매기기에서, 페이지랭크 \parencite{page1999pagerank}는 되풀이해서 퍼뜨린 보증으로 점에 순위를 매겨요. 이때 반복 한 번에 드는 비용은 화살표 수에 비례해서 늘어나요\parencite{langville2006pagerank}. 아이겐트러스트 \parencite{kamvar2003eigentrust}와 트러스트랭크 \parencite{gyongyi2004trustrank}는 정규화(합이 1이 되게 맞추기)와 씨앗 집합을 써서, 페이지랭크를 피어 투 피어 망의 신뢰와 웹 스팸에 맞게 바꾸었어요. 온체인 평판에서, \textcite{do2023hodgerank}의 연구는 이더리움 거래에 페이지랭크와 호지랭크를 적용했어요. AWP \parencite{do2023awp}는 거래 화살표에 금액과 나이에 따라 무게를 달아요. 그리고 자기 순위를 들어오는 송금의 개수와 금액과 견주어 점검해요. 이 통계는 AWP의 그래프와 같은 자료에서 나와요. \textcite{lin2025riskprop}에서는 계정의 위험을 거래 네트워크를 따라 퍼뜨려요. 블록체인 믿음·평판 시스템을 살펴본 조사 \parencite{bellini2020trust}와 스마트 계약 평판 설계를 살펴본 조사 \parencite{almasoud2020reputation}에 따르면, 흔히 쓰는 믿음 신호는 직접 매긴 평가와 송금량이에요. 이와 달리, 휴먼 패스포트 같은 이미 쓰이는 디파이 도구는 온체인과 오프체인의 자격 증명을 모아서, 한 주소가 서로 다른 한 사람의 것인지 정해요\parencite{humanpassportnddocs}. 그리고 코드 분석은 계약이 담고 있는 함수로 그 계약을 판단해요\parencite{lin2024crpwarner}. 이 연구들 가운데 어느 것도 ERC-20 허락 \parencite{vogelsteller2015erc20}을 순위 매기기 화살표로 쓰지 않아요. 또 계약이 받은 허락으로 계약에 순위를 매기지도 않아요. 주소 묶기(같은 주체의 주소를 한데 모으기)는 허락을 한 주체의 주소들을 잇는 데에만 써요\parencite{victor2020clustering}. 그리고 허락을 이용한 사기 수법을 분석한 연구는 속임수로 허락을 얻어 낼 수 있다는 것을 보여 줘요\parencite{adamczyk2025approval}. 디파이의 위험과 측정에서, 지식을 체계적으로 정리한 연구들은 위험을 서비스의 층에 따라 \parencite{werner2022sok}, 그리고 실제로 일어난 공격에 따라 \parencite{zhou2023sok} 나누어요. 자전 거래(자기끼리 사고팔아 거래가 많은 척하기) \parencite{victor2021washtrading}, 추출 가능한 가치(거래 순서를 이용해 빼내는 이익) \parencite{qin2022bev}, 사기 토큰 \parencite{xia2021scamtokens}, 에어드롭 파밍(공짜 토큰을 노린 가짜 활동) \parencite{luo2025airdrop}에 대한 연구는 송금 활동의 많은 부분이 꾸며 낸 것임을 보여 줘요. 측정 이론에서, 구성 타당도(점수가 재려는 것을 정말 재는지) \parencite{cronbach1955construct}는 점수가 재는 것과 바깥 기준을 구분해요. 그리고 부트스트랩 \parencite{efron1979bootstrap}은 순위 통계에 대해, 분포를 미리 가정하지 않는 구간을 줘요. 비어 있는 부분(\ref{sec:research-gap}절)은 허락 화살표로 계약에 순위를 매기고, 이것을 한 코호트에서 송금 방법 및 원시 차수와 나란히 평가하는 연구예요. 이 평가에는 신뢰구간, 짝지은 차이, 미리 정한 시간 재기 절차, 시간 홀드아웃이 들어가요. 시간 홀드아웃 가운데 하나는 등록 기간(자료를 뽑기 전에 등록해 둔 기간)에서 해요.

\dissertationSubheading[sec:methodology-summary]{연구 방법 요약}

\ref{ch:methodology}장은 공개된 장부 자료에, 수로 재고 서로 견주는 연구 설계를 써요. 분석 계획은 이 연구의 로그를 하나도 뽑기 전인 2026년 10월 9일에 커밋했어요. 이 계획은 코호트, 점수, 라벨, 통계, 규칙 A를 정해요. 2023년 10월 1일부터 2026년 9월 30일까지 아비트럼 원의 ERC-20 \texttt{Approval} 로그와 \texttt{Transfer} 로그는 공개된 구글 빅쿼리(BigQuery) 데이터 세트에서 뽑아요. 뽑는 범위는 계획의 코호트의 에고 네트워크(코호트 계약에 닿는 화살표만 모은 그래프)예요. 이 코호트의 규칙과 계정 종류 점검은 \ref{sec:data-sources}절에서 정의해요. 2026년 10월 10일의 수정은 등록 분석의 결과를 안 뒤에 했어요. 이 수정은 다섯 토큰으로 코호트를 다시 만들어요. 뽑은 로그 가운데 그 토큰들의 로그이면서 수정 코호트가 낀 것만 남겨요. 그 금액은 DefiLlama 서비스의 하루 가격으로 매겨요\parencite{defillamanddocs}. 그리고 금액 무게의 기울기를 중앙값 금액으로 정해요(\ref{sec:amount-revision}절). 주된 분석은 수정 분석이고, 등록 분석은 부록~\ref{sec:map-registered}에 남겨 두었어요. 거듭제곱 반복(같은 계산을 되풀이하기)으로 푸는 페이지랭크 구현 하나가 마지막 허락 그래프와 송금 그래프에 점수를 매겨요. 결합 연산자도 같은 풀이 프로그램 위에서 돌아가요. 각 점수는 스피어먼의 로($\rho$)와 켄달의 타우($\tau$)로 견주어요. 같은 기간에서는 송금, 허락, 시빌 안정성 대리 지표와 견주어요. 기간 밖에서는 W0과 W1의 스펜더 계약에서 새 허락 쌍, 새 송금자와 견주어요. 이때 각 동결일의 원시 차수를 기준선으로 써요. 모든 계수에는 계약을 짝지어 다시 뽑은 표본 400개로 구한 95\% 구간이 있어요(시드 42). 방법끼리의 비교는 짝지어서 해요. W1의 등록은 규칙 B를 더해요. 그리고 GMX~V2 계약 트레이더가 나중에 청산을 피하는지를 두고, 허락 쪽 점수와 송금 쪽 점수를 견주는 비교를 두 기간 모두에서 더해요. 이 라벨은 어느 화살표 종류로도 만들지 않았어요. 이 비교에서는 짝지어 다시 뽑은 표본 2{,}000개를 써요. 두 판정 대비에는 본페로니 97.5\% 구간(여러 번 비교할 때 넓힌 신뢰구간)을 써요. 고친 점수에는 규칙 B와 이 비교를 모두 재분석으로 적용해요. 실행 시간, 메모리, 반복 횟수, 그래프 크기는 한 가지 시간 재기 절차에 따라 기록해요. 모든 표의 바탕이 되는 코드, 설정, 컴퓨터가 읽을 수 있는 요약은 공개 저장소에 있어요(\ref{sec:reproducibility}절). 그리고 이 연구는 공개된 자료만 써요(\ref{sec:ethics}절).

\dissertationSubheading[sec:contributions]{기여}

이 논문은 온체인 평판과 디파이 위험 연구에 다섯 가지 기여(이 연구가 새로 보탠 것)를 해요. 엔도스랭크와 AWP는 둘 다 무게를 단 페이지랭크를 한 가지 화살표 종류에 적용해요\parencite{page1999pagerank}. 둘은 각자 자기 재시작 분포(어디서 다시 출발하는지의 비율)를 써요. 이동 규칙(다음에 어느 점으로 갈지 정하는 규칙)을 바꾸는 것은 결합 연산자뿐이에요. 그래서 새로운 것은 다음과 같아요. 화살표 종류 하나와 그것이 순위를 매기는 대상, 평가 설계, 허락 걷기가 무엇을 계산하는지에 대한 수학적 설명, 두 화살표 종류가 합쳐지는지에 대한 시험, 그리고 어느 화살표 종류에도 들어 있지 않은 결과 값에서의 비교예요. 이 비교는 주된 기간의 로그를 뽑기 전에 등록한 규칙을 따라요.
\begin{itemize}
\item 계약에 순위를 매기는 화살표로서의 허락 화살표(엔도스랭크): 보증 그래프에서 화살표 $u \to v$는 스펜더 $v$가 주인 $u$에게서 받아 가진, 기간 안의 다섯 토큰의 마지막 ERC-20 허락을 나타내요. 이 허락은 \texttt{approve}나 \texttt{permit}으로 정해져요. 금액은 달러로 매기고, AWP가 송금에 무게를 다는 방식대로 무게를 달아요. 이 화살표는 돈을 낸 기록이 아니라, 쓸 권리를 맡긴 일을 기록해요. 그리고 화살표는 스펜더를 가리켜요. 그래서 이 점수는 주인들이 허락해 준 계약에 순위를 매겨요. 허락은 한 주체의 주소들을 묶는 데 쓰인 적이 있어요\parencite{victor2020clustering}. 하지만 살펴본 평판 방법 가운데 계약이 받은 허락으로 계약에 순위를 매기는 것은 없어요(\ref{sec:research-gap}절).
\item 모든 점수를 그 점수가 다듬는 수와 견주는 평가, 같은 기간과 두 홀드아웃 기간에서: 모든 점수는 한 계약 묶음에서 계산해요. 그리고 AWP와 함께 쓰인 절차에 따라, 똑같은 송금, 허락, 시빌 안정성 대리 지표로 점검해요\parencite{do2023awp}. 모든 계수에는 95\% 부트스트랩 구간을, 모든 대비에는 짝지은 구간을 구해요\parencite{efron1979bootstrap}. 그다음 라벨 기간 전에 동결한 점수를, 스펜더 계약이 그 기간에 얻는 새 허락 쌍과 새 송금자로 시험해요. 이때 페이지랭크가 다듬는 원시 차수도 나란히 시험해요. 이 시험은 W0에서 하고, W1에서 다시 해요. W1의 코호트 정의와 라벨은 그 로그를 뽑기 전에 등록했고, 수정 분석에서는 다섯 토큰에 적용해요. 이 설계는 같은 기간의 일치와 예측을 나누어요. 또 퍼뜨린 점수와 그 점에서 그냥 센 수를 나누어요. 차수 대비는 미리 정한 규칙 어디에도 들지 않아요. 고친 점수에서는 두 기간의 라벨을 모두 안 뒤에 계산했어요. 이 자료에서 두 걷기는 달랐어요. 송금 걷기는 두 기간 모두에서 나중의 송금자를 원시 송금 수보다 더 잘 줄 세웠어요. 고르게 다시 출발하는 허락 걷기는 어느 기간에서도 나중의 허락을 그냥 센 허락한 주인 수보다 더 잘 줄 세우지 못했어요(주장~C4).
\item 깊이 1에서 엔도스랭크의 닫힌 꼴(한 줄 식으로 바로 쓴 답)과 퍼뜨림의 증거: 고르게 다시 출발하고, 매달린 점(나가는 화살표가 없는 점)의 몫을 고르게 다시 나누면, 주인들이 스스로는 허락을 하나도 받지 않은 스펜더는 $c\,(1+d\,\delta(v))$ 점을 받아요. 여기서 $c$는 들어오는 화살표가 없는 점들의 공통 점수예요. $d$는 감쇠 계수예요. $\delta(v)$는 주인 기준으로 나눈 들어오는 무게예요. 이 값은 각 주인의 보증을 그 주인이 허락한 모든 스펜더에게 나누어요(\ref{sub:theory-closed-form}절). 코호트의 허락 그래프에서, 이 식은 깊이 1인 모든 스펜더에서 1\% 안으로 맞아요. 하지만 들어오는 이웃이 더 깊이 이어진 스펜더 가운데에서는 절반도 안 되는 스펜더에서만 맞아요. 그래서 이 그래프에는 여러 걸음에 걸친 퍼뜨림이 있어요. 하지만 그 퍼뜨림 때문에 엔도스랭크가 나중의 허락을 원시 수보다 더 잘 줄 세우게 되지는 않았어요.
\item 결합 연산자와 그 시험: C-PR은 허락 층과 송금 층을 점마다 정해진 무게 $\lambda$로 섞어요. 두 층 모두 상한이 있는 값 변환(금액이 아무리 커도 1을 넘지 않게 바꾸는 식) 없이 달러 금액으로 무게를 달고, 허락 하나하나에는 상한을 둬요. S-PR은 씨앗을 쓰는 빼 보기 실험이에요. 규칙 A는 로그를 하나도 뽑기 전에 분석 계획에서 $\lambda$와 함께 정했어요. 규칙 A는 C-PR이 허락 층만 걷는 것보다 못하지 않고, 송금 층만 걷는 것보다 낫기를 요구했어요. 규칙 B는 W1의 로그를 뽑기 전에 등록했어요. 규칙 B는 C-PR이 송금 층만 걷는 것보다 낫기를 요구했어요. 등록 분석의 결과를 안 뒤에 고친 점수에 적용하면, 두 규칙 모두 실패해요. 등록 분석에서도 그랬어요. 이 판정은 등록한 시험이 아니라 재분석이에요. 그래도 새 허락 쌍에서 송금 걷기보다 나은 점은 W1에서 다시 나타났어요(주장~C5). 빼 보기 실험은, 미리 정하지 않은 대비에서, 허락 신호가 어떻게 되는지 보여 줘요. 허락 화살표가 걷기를 직접 이끌 때는 허락 신호가 살아남아요. 허락 화살표가 걷기가 어디서 다시 출발할지만 정할 때는 그 신호가 사라져요.
\item 어느 화살표 종류로도 만들지 않은 결과 값에서의 비교, 등록한 규칙에 따라: 허락 쪽 점수와 송금 쪽 점수를, GMX~V2 계약 트레이더가 나중에 청산을 피하는지를 두고 서로 견주어요. 이 라벨은 어느 그래프에도 들어가지 않아요. 판정 대비, 그 구간, 결론마다 쓸 문구는 7월--9월 로그를 뽑기 전에 등록했어요. 고친 점수로, 주된 기간에서, 재분석으로 하면, 등록한 문구는 두 결론을 줘요. 하나는 나중에 청산을 피하는 데에서, 그냥 센 수의 수준에서는 허락이 송금보다 앞서지 않는다는 결론이에요. 다른 하나는 페이지랭크 버전들은 어느 쪽이 나은지 가려지지 않는다는 결론이에요. 이익을 낸 비율과 손실이 아닌 비율에서는 허락한 주인 수가 송금 수보다 낮아요(주장~C6). 등록 분석도 같은 결론에 이르렀어요(부록~\ref{sec:map-registered}). 트레이더 표본은 작아요.
\end{itemize}


\dissertationSubheading[sec:scope-limitations]{범위와 한계}

이 연구는 아비트럼 원을 다뤄요. 관측 기간은 2023년 10월 1일부터 2026년 6월 30일까지 33개월이에요. 그리고 등록 기간 W1이 2026년 7월 1일부터 9월 30일까지 석 달(한 분기)을 더해요. 주된 표본은 관측 기간에 서로 다른 주인 세 명 이상에게서 WETH, USDC, USDT0, ARB, WBTC 가운데 하나의 0이 아닌 허락을 받은 계약 $n=15{,}107$개예요. 그래서 결과는 이 몇 달 동안 아비트럼 원에 있는 이런 종류의 계약에 들어맞아요. 더 적은 주인에게 허락을 받은 계약이나, 다른 토큰으로만 허락을 받은 계약에는 들어맞지 않아요. 홀드아웃 표본은 각 동결일에 이 토큰들 가운데 하나의 0보다 큰 마지막 허락을 가진 스펜더 계약이에요(W0에서 $n=17{,}050$, W1에서 $n=18{,}152$). 계정 종류는 2026년 10월 9일에 세 시간 안으로 고정한 블록들에서 부른 \texttt{eth\_getCode}로 정해요(\ref{sec:data-sources}절). 그래서 주소는 각 이벤트가 일어난 때의 코드가 아니라, 그 주소를 읽은 블록에서의 코드로 나뉘어요. 한쪽 끝에 코호트 계약이 있는 이벤트만 뽑아요. 그래서 각 그래프는 코호트의 에고 네트워크예요. 화살표 하나보다 긴 길은, 그 길 위의 모든 화살표가 코호트 계약에 닿을 때만 보여요. 그래서 여러 걸음에 걸친 퍼뜨림은 잘린 그래프에서 잰 것이에요(\ref{sub:theory-closed-form}절). 다섯 토큰의 이벤트만 그래프에 들어가요. 기본 화폐인 이더(ether)는 WETH로만 들어가요. 표~\ref{tab:usd-tokens}는 빠진 다른 토큰들의 행 수를 세어 보여 줘요(부록~\ref{sec:data-tokens}). 어느 화살표 종류로도 만들지 않은 결과 값에서의 비교는 계약인 GMX~V2 계정을 써요. 이 몇 달 동안 그런 계정 가운데 거래한 것은 적었어요(W0에서 $n=88$, W1에서 $n=70$). 그래서 모든 트레이더 결과는 구간이 넓어요(\ref{sub:neutral-label-results}절). 규모에 따른 실행 시간 변화는 코호트의 일부 표본에서 재요. 그리고 코호트의 에고 네트워크보다 큰 그래프는 시험하지 않아요.

하루 토큰 가격, 그리고 체인링크(Chainlink)의 가격 피드(블록체인에 올라오는 가격 정보) 목록 \parencite{chainlinknddatafeeds}에서 가져온 토큰 목록 말고는, 점수는 오프체인 정보를 쓰지 않아요. 보안 감사 보고서도, 사람이 골라 만든 계약 목록도, 휴먼 패스포트가 모으는 것 같은 신원 자격 증명도 쓰지 않아요. 어떤 점수도 해킹 공격, 러그 풀, 피싱을 찾아내는 도구로 시험하지 않았어요. 대출을 갚지 못한 일과 견주어 시험하지도 않았어요(\ref{sub:non-claims}절). 그리고 결과가 아비트럼 원이나 살펴본 몇 달을 넘어서도 들어맞는다고 주장하지 않아요. 금액은 그 토큰의 하루 가격으로 바꾼 달러예요. 이것은 등록 분석의 결과를 안 뒤에 한 선택이에요(\ref{sec:amount-revision}절). 엔도스랭크와 AWP는 금액을 상한이 있는 함수에 넣어요. 이 함수는 중앙값 금액에 절반의 무게를 줘요. 그래서 더 큰 허락이나 송금은 천장(얻을 수 있는 가장 큰 값)까지는 더 크게 세어져요. 무제한 허락은 다른 큰 허락과 똑같이 세어져요. C-PR의 층들은 금액 그 자체를 써요. 다만 허락 하나하나에 상한을 두어서, 무제한 허락이 그 허락이 들어 있는 주인의 줄(행)을 다 차지하지 않게 해요(\ref{ch:methodology}장). 허락이나 송금을 나쁜 뜻으로 조작하는 것은 실제로 있는 위협이에요. 주소를 만드는 데 돈이 거의 들지 않기 때문이에요\parencite{douceur2002sybil}. \ref{ch:discussion}장은 작은 시빌 주인 주소 농장(공격자가 만든 가짜 주소 무리)이 적은 비용으로 한 계약의 엔도스랭크를 코호트 대부분보다 위로 올릴 수 있다는 것을 보여 줘요. 반면에 코호트 꼭대기에 오르는 것은 비싸요(\ref{sec:sybil-cost-model}절). 토큰을 다섯 개만 남겨도 이 비용은 오르지 않아요. 허락에는 그 토큰의 잔고가 필요 없기 때문이에요. 시빌 보정 대리 지표는 몇몇 조작 방식을 보여 주지만, 점수를 그런 조작에서 지켜 주지는 않아요.

\dissertationSubheading[sec:roadmap]{연구 길잡이}

\ref{ch:literature}장은 그래프에 바탕을 둔 온체인 평판, ERC-20 허락의 뜻, 디파이에서 쓰이고 있는 평판 도구를 살펴봐요. 그리고 이론의 틀을 세워요. \ref{ch:methodology}장은 코호트, 토큰, 가격과 고친 금액, 그래프와 점수, 구성개념 점검, 시간 홀드아웃 W0과 W1, 시간 재기 절차, 그리고 재현성(다른 사람이 똑같이 다시 해 볼 수 있음)과 윤리의 기준을 정의해요. \ref{ch:results}장은 수정 분석을 보고해요. 벤치마크(속도 재기, 주장~C1), 모든 점수의 같은 기간 일치도(C2), 순위 갈라짐(C3), 감쇠 계수, 금액 기울기, 표본 정의, 다시 출발에 대한 점검, 등록 기간을 포함한 시간 홀드아웃(C4와 C5), 그리고 어느 화살표 종류로도 만들지 않은 결과 값에서의 비교(C6)예요. 이때 등록 분석을 위해 정한 규칙을 재분석으로 적용해요. 부록~\ref{ch:appendix-er-awp}에는 엔도스랭크와 AWP를 직접 견준 결과를 모았어요. 부록~\ref{sec:map-registered}은 등록 분석을 수정 분석과 나란히 보여 줘요. \ref{ch:discussion}장은 찾아낸 것을 연구 질문과 목표에 따라 풀이해요. 그 뜻하는 바를 따져 보고, 허락 그래프에 대한 시빌 공격의 비용을 모형으로 만들어요. 그리고 타당성을 위협하는 것과 한계를 따져 봐요. \ref{ch:conclusion}장은 RQ1--RQ5에 답하고, 지식에 대한 기여를 밝히고, 앞으로 할 일을 내놓아요.
